% Complete research fragment. The finite coefficient proof is reproduced below.

\section{An exact weight-five normalization}
\label{sec:weight-five-normalization}

Let $G=\widehat{\mathrm{GT}}$, let $H$ be its Galois image, and write
the radical and cubic normalizations as
\[
 D=N_{\mathrm{rad}},\qquad
 D_3=D\cap\ker m_3,\qquad J_3=D_3\cap H.
\]
The established factorizations give $G=D_3H$.  All Magnus coefficients
in this section are $2$-adic.  For $g=(1,f)\in D_3$, the pro-$2$
Magnus expansion of $f$ starts in degree five: the degree-three and
degree-four Hall coordinates are respectively
$(m_3,-m_3)$ and $(-m_3,m_3,-m_3)$ by the preceding low-weight
calculation.  Define
\[
 m_5(g)=[S^4T]f,\qquad x=1+S,\quad y=1+T.
\]

\begin{theorem}\label{thm:weight-five-exact-image}
The exact integral images are
\[
 m_5(D_3)=m_5(J_3)=32\mathbf Z_2.
\]
In particular,
\[
 \eta_5:D_3\longrightarrow\mathbf F_2,
 \qquad \eta_5(g)=m_5(g)/32\pmod2,
\]
is a continuous $G$-invariant character with nonzero arithmetic
restriction.  The kernel of $m_5$,
\[
 D_5=D_3\cap\ker m_5,
 \qquad J_5=D_5\cap H=J_3\cap\ker m_5,
\]
is a proper closed normal subgroup satisfying
\[
 D_3=D_5J_3,\qquad G=D_5H.
\]
The relative abelian quotient refines canonically to
\[
 B\simeq D_3/\overline{[G,D_3]J_3}
       \simeq D_5/\overline{[G,D_5]J_5}.
\]
\end{theorem}

We first prove the upper bound.  The only global GT input needed
for this step is the normalized signed congruence identity
\begin{equation}\label{eq:weight-five-signed-input}
 f(U^2,V^2)=I,
 \qquad
 U=\begin{pmatrix}1&1\\0&1\end{pmatrix},\quad
 V=\begin{pmatrix}1&0\\-1&1\end{pmatrix}.
\end{equation}
The remaining input consists of the two-cycle and three-cycle
relations in the free pro-$2$ group.

\begin{lemma}[Finite integral coefficient relations]
\label{lem:weight-five-coefficients}
Suppose $f\in\gamma_5$ satisfies
\[
 f(y,x)f(x,y)=1,\qquad
 f(z,x)f(y,z)f(x,y)=1,\qquad z=(xy)^{-1}.
\]
Put
\[
 a=[S^4T]f,\quad
 c_5=[STSTS]f,\quad
 c_7=[STSTSTS]f,\quad
 c_9=[STSTSTSTS]f.
\]
Then
\begin{equation}\label{eq:weight-five-three-coefficients}
 c_5=-2a,\qquad c_7\equiv5a\pmod{16},\qquad
 c_9\equiv2a\pmod4.
\end{equation}
\end{lemma}

The explicit finite recursions and coefficient certificates proving
this lemma are given below.  In particular the third relation does
not omit a product of lower-weight terms: products of two
nonconstant factors begin in degree ten.

\begin{proposition}\label{prop:weight-five-upper-bound}
For $g\in D_3$, one has $m_5(g)\in32\mathbf Z_2$.
\end{proposition}

\begin{proof}
At the matrices in \eqref{eq:weight-five-signed-input}, consecutive
equal Magnus letters vanish.  Surviving even-length words are
diagonal.  The upper-right entry through degree nine is
\[
 32c_5-128c_7+512c_9.
\]
Write $c_7=5a+16r$.  Modulo $512$, the matrix identity gives
\[
 0\equiv32(-2a)-128(5a)=-64\cdot11a\pmod{512},
\]
so $a\in8\mathbf Z_2$.  The third relation in
\eqref{eq:weight-five-three-coefficients} now gives $c_9\in4\mathbf Z_2$.
Degree ten is diagonal, and every term of degree at least eleven
is divisible by $2048$.  Reducing the same entry modulo $2048$
therefore gives
\[
 0\equiv-64\cdot11a\pmod{2048}.
\]
Since $11$ is a $2$-adic unit, $a\in32\mathbf Z_2$.
\end{proof}

\subsection{Finite coefficient certificate}
\label{subsec:weight-five-certificate}

For a word $w$ in the alphabet $x<y$, write $[w]$ for its Magnus
coefficient after replacing $x,y$ by $S,T$.  A word is Lyndon if
it is lexicographically smaller than every nontrivial cyclic
rotation.  Define the standard Lyndon group words by
\[
 L_x=x,\quad L_y=y,\qquad L_w=[L_u,L_v]
\]
when $w=uv$ is the standard factorization with $v$ the longest
proper Lyndon suffix.  Let $h_w\in\mathbf Z_2$ be the corresponding
Hall coordinates of $f$.  Up to degree nine,
\begin{equation}\label{eq:weight-five-linear-collection}
 f=1+\sum_{5\le |w|\le9}h_w(L_w-1).
\end{equation}
This is an equality of truncated Magnus expansions; powers and
products have no quadratic contribution before degree ten.

Here are the complete recursions used to check the coefficient
certificates.  In the ring of noncommutative integer polynomials
truncated after degree eight, put
\[
 X=1+S,\quad Y=1+T,\quad Z=Y^{-1}X^{-1},\qquad
 (1+A)^{-1}=\sum_{j=0}^{8}(-A)^j.
\]
Form $L_w$ by the preceding commutator recursion and put
$F=1+\sum_{5\le|w|\le8}h_w(L_w-1)$.  Define the linear coefficient
forms
\[
 I_v=[v]\bigl(F(Y,X)+F(X,Y)-2\bigr),
\]
\[
 II_v=[v]\bigl(F(Z,X)+F(Y,Z)+F(X,Y)-3\bigr).
\]
All $I_v$ and $II_v$ vanish for an $f$ in the lemma.  The following
identities are identities of integer linear forms in the Hall
coordinates, with the indicated modulus; thus they may be checked
by the displayed finite multiplication and substitution recursions.

First,
\begin{align*}
 3(c_5+2a)={}&36I_{xxxxy}-18I_{xxxyy}+24I_{xxyxy}
             +66I_{xxxxyy}+24I_{xxxyxy}+6I_{xxyyxy}\\
             &+26II_{xxxyy}+66II_{xxxxxy}-12II_{xxxyxy}.
\end{align*}
This proves $c_5=-2a$ without division by a nonunit.

Second, modulo $16$ one has
\begin{equation}\label{eq:weight-five-c7-certificate}
 9(c_7-5a)\equiv\sum_v\alpha_v I_v+\sum_v\beta_v II_v,
\end{equation}
where the complete lists of nonzero coefficients are
\[
\begin{array}{c|r@{\qquad}c|r}
 v&\alpha_v&v&\beta_v\\ \hline
 xxxxy&-4&xxxyy&-5\\
 xxxyy&-1&xxxxxy&-7\\
 xxyxy&6&xxxyxy&-6\\
 xxxxyy&5&xxyxyy&8\\
 xxxyxy&-4&xxxxxyy&4\\
 xxyyxy&-1&xxxyxxy&4\\
 xxxxyyy&8&xxxyxyy&8\\
 xxxyxyy&8&xxxyyxy&8\\
 xxyxxyy&-4&xxxxxyyy&8\\
 xxyxyxy&8&xxxxyxyy&8\\
 xxxyyxyy&8&xxxyxyxy&8\\
 xxyxyyxy&8&xxyxxyxy&8.
\end{array}
\]
Since $9$ is a unit modulo $16$, this proves the second assertion
of \eqref{eq:weight-five-three-coefficients}.

For the degree-nine assertion, the only Hall coordinates that
can contribute modulo four are determined by a smaller matrix
recursion.  Work in $\mathbf Z[t]/(t^{10})$ and put
\[
 P_x=\begin{pmatrix}1&t\\0&1\end{pmatrix},\qquad
 P_y=\begin{pmatrix}1&0\\-t&1\end{pmatrix},\qquad
 P_w=P_uP_vP_u^{-1}P_v^{-1}.
\]
The inverse is the usual determinant-one matrix inverse, followed
by truncation.  Then $d_w=[t^9](P_w)_{12}$ is precisely
$[STSTSTSTS]L_w$.  The full table relevant modulo four is
\[
\begin{array}{c|r|r|l}
 |w|&\text{number of Lyndon words}&\gcd_w d_w&
       \{w:d_w\not\equiv0\pmod4\}\\ \hline
 5&6&2&xxxxy,\ xxxyy,\ xxyxy\quad(d_w\equiv2)\\
 6&9&2&xxxyxy\quad(d_w\equiv2)\\
 7&18&4&\varnothing\\
 8&30&8&\varnothing\\
 9&56&16&\varnothing.
\end{array}
\]
For additional explicit checks, the unabridged degree-nine
coefficient vectors in weights five and six, in lexicographic
Lyndon order, are
\[
 (-2,-22,18,4,0,0),\qquad
 (0,-8,14,4,8,-4,0,0,0).
\]
The recursion specifies all the remaining entries; their
divisibility by four is the only property used here.

There is also the following four-row identity of Hall coefficient
forms modulo two:
\begin{equation}\label{eq:weight-five-c9-parity-certificate}
 h_{xxxyy}+h_{xxyxy}+h_{xxxyxy}
 \equiv I_{xxxxyyy}+I_{xyxxxyyy}+II_{xxyxxy}+II_{xyxyyyy}
 \pmod2.
\end{equation}
Since $h_{xxxxy}=a$, the table, linear collection
\eqref{eq:weight-five-linear-collection}, and
\eqref{eq:weight-five-c9-parity-certificate} give
\[
 c_9\equiv
 2(h_{xxxxy}+h_{xxxyy}+h_{xxyxy}+h_{xxxyxy})
 \equiv2a\pmod4.
\]
This completes the finite certificate and the proof of
Lemma~\ref{lem:weight-five-coefficients}.

\subsection{An arithmetic generator of the exact image}

The arithmetic lower bound uses a weight-five homomorphism on
the entire multiplier-one subgroup.  For such an element, write
\[
 a_3=[S^2T]f,\qquad b_5=[S^4T]f,
 \qquad \tau_5(f)=12b_5-11a_3\in\mathbf Z_2.
\]

\begin{lemma}\label{lem:tau-five}
The map $\tau_5$ is a continuous homomorphism on $\ker\lambda$,
and
\[
 \tau_5(hzh^{-1})=\lambda_2(h)^5\tau_5(z)
 \qquad(h\in G,\ z\in\ker\lambda).
\]
On $D_3$, it is equal to $12m_5$.
\end{lemma}

\begin{proof}
Let $R(X,Y)$ be the Alexander coefficient of $f$ on $[x,y]$,
and let
\[
 B_f(X,Y)=1+(X-1)(Y-1)R(X,Y).
\]
The action on the rank-one Alexander module is semilinear.
Consequently, on the multiplier-one subgroup the cofactors
multiply, and conjugation by multiplier $t$ sends $B_f(X,Y)$
to $B_f(X^t,Y^t)$.  It follows that
\[
 [u^4v]\log B_f(e^u,e^v)
\]
is an additive weight-five $\mathbf Q_2$-valued homomorphism.
Only terms linear in $Y-1$ contribute to this coefficient, so
the logarithm can be omitted for this extraction.  The boundary
Magnus coefficients and the earlier identity $[S^3T]f=-a_3$ give
\[
 (X-1)R(X,1)
 =a_3(X-1)^2-a_3(X-1)^3+b_5(X-1)^4+O((X-1)^5).
\]
Now
\[
 [u^4](e^u-1)^2=\frac7{12},\quad
 [u^4](e^u-1)^3=\frac32,\quad
 [u^4](e^u-1)^4=1.
\]
Thus the extracted homomorphism is $b_5-11a_3/12$.
Multiplication by $12$ proves the asserted integral formula
and its exact transformation rule.  The last assertion uses
$a_3=0$ on $D_3$.
\end{proof}

\begin{proposition}\label{prop:weight-five-arithmetic-generator}
There exists $w\in J_3$ with $m_5(w)\in32\mathbf Z_2^\times$.
\end{proposition}

\begin{proof}
Simultaneous rational Kummer surjectivity over
$\mathbf Q^{\mathrm{ab}}$ supplies $\sigma\in H\cap\ker\lambda$
with
\[
 \rho_2^I(\sigma)=2\in\widehat{\mathbf Z},\qquad
 \rho_q^I(\sigma)=0\quad(q\text{ odd}).
\]
The established boundary congruence
$\rho_2^I\equiv-2m_3\pmod4$ makes $a_3(\sigma)$ odd.
Therefore $\tau_5(\sigma)=12b_5(\sigma)-11a_3(\sigma)$ is odd.

Choose $h\in H$ with $\lambda_2(h)=3$ and put
$t=\lambda(h)\in\widehat{\mathbf Z}^{\times}$.  For
$c\in\widehat{\mathbf Z}$ and $z\in H\cap\ker\lambda$, define
\[
 \mathcal D_c(z)=hzh^{-1}z^{-c}.
\]
The power $z^{-c}$ is the standard continuous profinite power
in the procyclic closure of $z$.

Every rational Kummer coordinate of $\mathcal D_t(\sigma)$
is zero, because its conjugation weight is $t$.  Set
\[
 w=\mathcal D_{t^3}\bigl(\mathcal D_t(\sigma)\bigr).
\]
All rational Kummer coordinates remain zero.  Its cubic
coordinate also vanishes, since conjugation on that coordinate
has weight $t_2^3$.  Hence $w\in J_3$.

Lemma~\ref{lem:tau-five} and the continuity of powers give the
exact equality
\[
 \tau_5(w)=(3^5-3^3)(3^5-3)\tau_5(\sigma)
          =51840\tau_5(\sigma).
\]
This equality holds for the exact homomorphism $\tau_5$ with the
stated conjugation weight.  Since $w\in D_3$,
\[
 m_5(w)=\frac{51840}{12}\tau_5(\sigma)
       =4320\tau_5(\sigma)
       =32\cdot135\,\tau_5(\sigma),
\]
which is a generator of $32\mathbf Z_2$.
\end{proof}

\begin{proof}[Proof of Theorem~\ref{thm:weight-five-exact-image}]
On $D_3$, $m_5$ is a continuous additive homomorphism, with
conjugation weight $\lambda_2^5$; this follows either from its
leading Magnus degree or from Lemma~\ref{lem:tau-five} and
the torsion-freeness of $\mathbf Z_2$.  The upper bound and
the arithmetic generator show that its compact images on
both $D_3$ and $J_3$ are exactly $32\mathbf Z_2$.  Dividing
by $32$ and reducing modulo two gives a $G$-invariant
character, since every $2$-adic unit has odd fifth power.

Equality of the images gives $D_3=D_5J_3$.  Normality follows
from the scalar conjugation rule, properness from the
arithmetic generator, and $G=D_5H$ follows from $G=D_3H$.
For the quotient assertion, if $G=EH$ with $E$ closed normal,
then
\[
 E/\overline{[G,E](E\cap H)}
 \simeq G/\overline{[G,E]H}.
\]
If $E=E'(E\cap H)$ and $G=E'H$, the closures
$\overline{[G,E]H}$ and $\overline{[G,E']H}$ coincide: in
the quotient by the latter, $H$ is trivial and the image
of $E'$ is central.  Apply this with $E=D_3$ and $E'=D_5$.
\end{proof}

If $J_3=G_{L_3}$ and $L_5=\overline{\mathbf Q}^{\,J_5}$,
then $\operatorname{Gal}(L_5/L_3)\simeq\mathbf Z_2$.
The character $\eta_5$ has nonzero arithmetic restriction on $D_3$.
The same witness shows precisely why it does not give another invariant
character on the original normalized subgroup.

\begin{corollary}\label{cor:weight-five-nonextension}
The character $\eta_5$ does not extend to a $G$-invariant continuous
$\mathbf F_2$-character of $D=N_{\mathrm{rad}}$, and hence does not
extend to one of $N$.  Nevertheless $D_5$ is a strictly smaller
normalization with the same relative abelian quotient.  The kernel of
restriction from $X_2$ to $H^1(D_5,\mathbf F_2)$ is exactly
\[
 \bigoplus_{q\text{ odd prime}}\mathbf F_2\psi_q
       \oplus\mathbf F_2\eta_3.
\]
The kernel of restriction from $H^1(D,\mathbf F_2)^G$ to
$H^1(D_5,\mathbf F_2)$ is exactly $\mathbf F_2\eta_3|_D$.
For every odd prime $p$, restriction from $H^1(N,\mathbf F_p)^G$
to $H^1(D_5,\mathbf F_p)$ is injective.
\end{corollary}

\begin{proof}
Use the elements $h,\sigma,t$ of
Proposition~\ref{prop:weight-five-arithmetic-generator}, and put
$z=\mathcal D_t(\sigma)\in D\cap H$.  Its witness is
\[
 w=hzh^{-1}z^{-t^3}=[h,z]z^{1-t^3}\in D_3\cap H,
 \qquad \eta_5(w)=1.
\]
Since $t$ is a profinite unit, $1-t^3\in2\widehat{\mathbf Z}$.
The continuous profinite power $z^{(1-t^3)/2}$ is defined, so
$w$ belongs to the closed subgroup generated by $[G,D]$ and the
squares in $D$.  Every invariant quadratic character of $D$ kills
that subgroup.  It therefore takes the value zero at $w$, which is
incompatible with extending $\eta_5$.  Restricting an invariant
character of $N$ to $D$ proves the second nonextension assertion.

If $\chi\in X_2$ vanishes on $D_5$, its restriction to $D_3$
factors through $D_3/D_5\simeq\mathbf Z_2$.  That restriction is
either zero or $\eta_5$.  The latter possibility has just been
excluded, so $\chi$ vanishes on $D_3$.  Conversely every character
vanishing on $D_3$ vanishes on $D_5$.  Corollary~\ref{cor:cubic-radical-joint-image}
identifies this common kernel as the displayed direct sum.
The same argument on $D$, using $D/D_3\simeq8\mathbf Z_2$,
gives the one-dimensional kernel there.

For odd $p$, a character of $D_3$ vanishing on $D_5$ must be zero,
because $D_3/D_5$ is pro-$2$.  A character of $N$ vanishing on
$D_3$ factors through the joint rational Kummer and cubic quotient.
The cubic factor is pro-$2$, and conjugation by multiplier $-1$
acts by inversion on every rational Kummer coordinate.  Invariance
therefore kills the remaining factor in $\mathbf F_p$.
\end{proof}
